\hypertarget{the-service-mesh}{%
\chapter{14. The service mesh}\label{the-service-mesh}}

A single kennel is complete on its own. It holds the capabilities its
policy grants, reaches nothing else, and two kennels are mutually
invisible by default. A mesh is the deliberate exception: a workload
that uses a capability another confined kennel provides, with neither
kennel holding the other's authority. A desktop application reaches a
display server it does not run; an application reaches a cache without
seeing the process behind it. The lateral model holds throughout,
provider and consumer siblings under the daemon and neither nested in
the other, and the only thing that crosses between them is the one
connector the operator declared and the daemon brokered.

One declaration carries the whole surface. A provider lists each
capability it offers under \texttt{{[}{[}provides{]}{]}}, by a public
name, a typed shape, and the \texttt{endpoint} where it exposes the
capability in its own view; a consumer lists each it reaches under
\texttt{{[}{[}consumes{]}{]}}, by the same name and shape and the
\texttt{at} where the broker presents it for the workload to act
against.

\begin{kennelcode}
# in the provider's policy
[[provides]]
name = "org.projectkennel.dbus-broker"
shape = "binder-connector"
endpoint = "/dev/binderfs-mesh/binder"
reason = "push per-session ACCEPT_SESSION to the broker's control node on the mesh bus"

# in a consumer's policy
[[consumes]]
name = "org.projectkennel.dbus"
shape = "dbus-name"
required = true
reason = "reach the standing D-Bus broker's mediation over the connector mesh"
\end{kennelcode}

Neither side names the other. A consumer names the capability the way a
client names a service, never the process behind it, and nothing in
either declaration points at another kennel, so there is no cross-kennel
reference anywhere in the policy surface to resolve at compile. The
shape names the transport, and the surface carries three:
\texttt{af-unix}, \texttt{dbus-name}, and \texttt{binder-connector}. The
\texttt{af-unix} shape the broker delivers directly and the display
service rides it; the other two ride the mesh bus, taken up where the
mesh is (§14.6). Everything before that holds for all three, because
declaration, catalogue, matching, and enablement are shape-agnostic and
only the final handoff is not.

\hypertarget{compile-local-resolve-at-runtime}{%
\section{14.1 Compile-local, resolve at
runtime}\label{compile-local-resolve-at-runtime}}

A cross-kennel match cannot be settled when a policy is compiled, and
the design does not pretend otherwise. Policies are compiled and signed
independently, so the compiler only ever holds the one policy in front
of it, and the providers that can answer a consume are whatever the
operator has installed by the time the workload reaches for it, not
whatever existed when the consumer was signed. A signed artefact cannot
be made to depend on the presence of another mutable, independently
authored one, and resolving a consume against the known kennels at
compile has no known kennels to resolve against. Matching is therefore a
runtime property, and putting it anywhere else would be a guarantee the
system could not keep.

So the two checks live in two places. At compile, the validation is
local: each declaration is well-formed, the shape is one of the defined
transports, a reserved name traces to a template signed by a key
authorised for it (§14.4), and no two \texttt{{[}{[}provides{]}{]}} in
the one policy claim the same name. None of these consults another
kennel, so all of them hold on a signed artefact for its whole life. At
runtime, the broker resolves a consumer's name against the catalogue,
requires the optional key to match, enforces the expected shape, and
applies the deny-by-default identity check. A consume marked
\texttt{required\ =\ true}, the default, has its name checked for
resolvability at construction, and a kennel whose hard dependency
nothing provides fails to start rather than run missing something it
needs (\texttt{refuse-to-start}); a \texttt{required\ =\ false} consume
is skipped when unavailable and the workload tolerates the absence, the
systemd \texttt{Requires=} and \texttt{Wants=} distinction evaluated
against the installed set. Tooling may report which of a deployment's
consumes currently resolve, but that is an operator diagnostic over a
set that changes the moment a service is added or removed, explicitly
advisory and never a signature gate.

\hypertarget{deny-by-default-and-the-broker}{%
\section{14.2 Deny-by-default and the
broker}\label{deny-by-default-and-the-broker}}

Authorisation is the consumer's signed \texttt{{[}{[}consumes{]}{]}} and
nothing else. A kennel reaches a capability if and only if its own
operator-signed policy declares a consume for it, and a kennel with no
such declaration reaches nothing. The provider is passive: it offered
the capability, and the operator decides who uses it by what it signs
into each consumer's policy, never by an allowlist the provider
maintains. This keeps the provider decoupled from its consumer set, the
display service needing no knowledge of which applications render
through it, and it is the same shape as the rest of the system, where a
capability is granted on the kennel that uses it exactly as egress is
granted on the consuming kennel and not as an allowlist on the
destination.

When a workload first acts against its \texttt{at}, the broker runs a
fixed sequence, each step a lookup or an equality rather than a fresh
policy evaluation. It checks the request against the consumer's
kernel-stamped identity and confirms the signed policy declares this
consume, so the workload can request only what its policy already grants
and cannot widen it (\texttt{request-dont-author}). It resolves the name
against the catalogue to a single provider. It requires the optional key
to match exactly, where the key is a discriminator both sides set to the
same opaque value, so a keyed consumer is never brokered to a keyless
provider and the public name a different kennel could also advertise
does not by itself land a consumer at the wrong one. It enforces the
expected shape, refusing a mismatched transport rather than
misdelivering. Then it activates the provider if it is cold, bridges the
connection, and steps out of the byte path. A failure at any step, no
installed provider, a shape that disagrees, a key that does not match,
is a denial and an audit record, never a silent fallback to some other
provider.

Two properties of the name keep the broker honest. A reserved
\texttt{org.projectkennel.*} name is claimable only by a template signed
with the project maintainer key, so a workload cannot advertise the
display service's name and have a consumer brokered to an impostor, the
provider-name spoof the catalogue would otherwise carry. And a public
name offered by more than one enabled provider is never collapsed to no
winner: collapsing would let one provider revoke a name another serves
merely by also claiming it, a denial of service by name-claim, so a
second claimant adds a candidate and can never empty the name, the
broker selecting among candidates by key and then by the
per-user-over-per-host cascade. Through all of it the daemon brokers the
connector and parses none of what rides it, which is what keeps it from
becoming the relay its centrality would otherwise make it
(\texttt{control-not-data-plane}, T5.3).

\hypertarget{the-catalogue-is-the-filesystem}{%
\section{14.3 The catalogue is the
filesystem}\label{the-catalogue-is-the-filesystem}}

The catalogue the broker resolves against is not standing daemon state.
It is a derived projection of the signed \texttt{{[}{[}provides{]}{]}}
blocks of the providers the operator has enabled, and enabled means
linked: the operator turns on what a vendor provides by placing a
symlink to its signed policy into an enablement directory, eager in
\texttt{autorun/} or lazy in \texttt{ondemand/}. The split between what
the signature covers and what the link covers is deliberate. The signed
provider policy carries the capability and the supervision discipline,
what the service is and how it tolerates its own crashes, which is the
author's to vouch for; whether it autostarts, and whether eagerly or
lazily, is the operator's deployment posture, expressed solely by which
directory holds the link and kept out of the signed artefact. A vendor
cannot bake autostart-on-every-host into a signed policy, and an
operator cannot alter a service's restart discipline without re-signing
it. Each side owns exactly its half.

Because the enabled set is the links on disk,
\texttt{kennel\ daemon-reload} re-scans the enablement directories,
re-derives the catalogue, brings newly linked eager providers online,
and makes newly linked lazy providers resolvable without starting them;
removing a link and reloading drops the capability and stops an eager
provider. The filesystem is the registry, re-read on every reload and on
daemon restart, so the daemon holds no enabled-set state a restart could
lose or a bug could desynchronise.

The lazy posture is what makes a capability reachable before its
provider runs. A provider linked into \texttt{ondemand/} is
socket-activated on the first consume, the daemon doing for a provider
kennel what systemd's socket activation does for the daemon itself,
while an eager provider is simply already up. This is why delivery is a
socket the workload connects to and not a connected descriptor handed in
at the workload's own construction: an inherited descriptor would have
to exist before any consumer reached for it, which forces the provider
up eagerly and defeats the lazy model, whereas a socket standing in the
consumer's view from construction costs nothing while idle and makes the
connect itself the trigger that resolves and activates the provider. A
connect against a cold provider therefore blocks until the provider is
ready or a deadline fires, a broker constant defaulting to five seconds,
and a provider still pending does not satisfy a waiter. That single rule
makes the flat, boot-order-free model safe under a dependency cycle: two
services that consume each other, both enabled eager, each block on the
other's readiness, both hit the deadline, and both resolve to a loud,
observable double-timeout visible in the topology surface rather than a
silent deadlock. A brokered connector is a live handle and not a
standing guarantee, so when a provider restarts or dies its consumers
observe an end-of-file and re-request the capability, which re-resolves
and activates afresh; the mesh guarantees a connection to the named
capability on request, never the same connection for the kennel's life.

\hypertarget{supervision-and-the-service-kennel-trust-class}{%
\section{14.4 Supervision and the service-kennel trust
class}\label{supervision-and-the-service-kennel-trust-class}}

A provider the operator enables is a kennel the daemon keeps running on
the operator's behalf, so it carries the one thing an ephemeral spawn
does not, a supervision discipline declared in a \texttt{{[}service{]}}
block and signed into the policy. The block names a restart policy, a
backoff that doubles on each successive attempt so an immediately
crashing provider does not spin the supervisor, and a bound on attempts
within the crash-loop window. An enabled provider is in one of three
readiness states the catalogue projects and the topology surface reads:
declared-but-pending, enabled and known but not yet constructed;
declared-and-ready, reachable, a consumer's connect bridging straight
through; and declared-but-failed, where construction or supervision gave
up. A failed provider stays in the catalogue, so its failure is visible
and a consume against it is denied and distinguishable from
no-such-capability, and it is terminal until the operator intervenes, a
failure made sticky and observable rather than self-clearing.

Who may claim a name is the load-bearing gate, because the name is the
thing a consumer trusts: a kennel that resolves
\texttt{org.projectkennel.wayland} is trusting whatever the mesh brokers
it under that name to be the display service. The
\texttt{org.projectkennel.*} namespace is the project's own, built in
and claimable only by a maintainer-signed template, the same unit of
trust a spawn target uses; an unreserved name is freely authored and
signed by any valid key. A maintainer-signed reserved template is not
frozen whole, exposing the bits a host legitimately varies as the same
fenced mutable fields a spawn target declares, so a host runs a
different inner compositor without re-authoring the reserved name or
escaping its trust. A host may additionally reserve its own namespaces
in the root-owned deployment config, purely additive and never able to
redefine the project's.

The service-kennel trust class carries one deliberate widening of the
single-leg discipline. The rule that an agent-composable spawn target
holds at most one leg binds what an untrusted agent composes, because
holding two at once would let it reconstitute across one kennel a
capability the operator never granted whole. A service kennel is not
that: it is a maintainer-signed template the operator deliberately
enabled, vouched for by both the signature and the enablement, so it may
hold several legs, the GUI service the worked member of the class. The
exemption widens how many legs a service kennel holds, never what kind.
It may vend capabilities a consumer uses, render, transport, reach a
destination, but never capabilities whose purpose is to be trusted, a
secrets broker or a signing service, which would place a trust root
inside the boundary the project exists to confine
(\texttt{authentication-never-attestation}). Trust material a kennel
needs arrives as a signed construction parameter from the operator,
never brokered from a peer at runtime. The consume grant is loud, each
declaration carrying a reason surfaced in the risk report, and the
standing shape leaves two narrowed residuals the catalogue records: a
provider is a shared target whose compromise reaches every consumer it
serves for its lifetime, and the brokering edge is a standing addition
to the daemon's always-present surface, bounded by the parse-nothing
discipline rather than removed (T3.10).

\hypertarget{the-af-unix-shape}{%
\section{14.5 The af-unix shape}\label{the-af-unix-shape}}

The delivered shape is \texttt{af-unix}, and its handoff is the same
fd-broker every other capability rides. On a successful consume the
broker attaches a connected descriptor to the provider, the byte path
the daemon then steps out of. Where that descriptor connects is a
host-owned rendezvous point: the daemon derives a per-capability
directory under its runtime root from the signed catalogue triple of
tier, name, and optional key, binds that directory into the provider's
view at the parent of its declared endpoint, and connects to the
host-side inode the provider binds its socket on, a bind mount being one
inode set seen from two namespaces. This is byte-for-byte the
host-socket facade that brokers an ordinary
\texttt{{[}{[}unix.allow{]}{]}} socket, pointed at a provider's endpoint
instead of a host path, and it binds a provider's own capability
directory alone rather than a shared parent, so a provider sees only its
own rendezvous socket and every cross-kennel reach stays brokered and
deny-by-default. The confined display service is the first non-trivial
consumer of all of this, its host-compositor leg riding this shape.

\hypertarget{the-mesh-bus}{%
\section{14.6 The mesh bus}\label{the-mesh-bus}}

A capability that crosses between two kennels needs a binder instance
both can reach, and a kennel's own bus is private to it, so the mesh
shapes ride a different one. The mesh bus is a single binder instance
distinct from every per-kennel bus, with \texttt{kenneld} on its node 0,
on which a provider registers a node by \texttt{ADD\_SERVICE} and a
consumer receives a handle to that node by \texttt{SVC\_CONNECT}. The
per-kennel bus carries a kennel's calls to the daemon; the mesh bus
carries one kennel's reach to another, and the daemon holds node 0 of
both.

The instance exists where an unprivileged daemon cannot ordinarily put
it. Mounting a binderfs needs the capabilities a user namespace confers,
which \texttt{kenneld} holds only inside one it creates, so the mesh
binderfs is mounted by a child the daemon forks for the purpose. The
child stays under the daemon's own AppArmor profile, which carries the
user-namespace grant across a fork and not across an exec, so the fork
is load-bearing where an exec would lose the grant. It maps itself to
the daemon's uid in a single-uid namespace, unprivileged and with no
subordinate uids, mounts the binderfs, and the nodes it creates are
owned by the daemon's own uid. \texttt{kenneld} serves node 0 by opening
the device through the holder's
\texttt{/proc/\textless{}holder\textgreater{}/root}, the same reach by
which it holds node 0 of every per-kennel bus, pointed at the holder
instead of an init. No privileged helper takes part: the mount the
privilege model refuses the daemon directly is the one it makes for
itself inside a namespace it is sovereign over.

Exposing the device into a consumer's view is a second problem, because
the mount lives in the holder's mount namespace and a kennel's
construction has a private pid namespace in which
\texttt{/proc/\textless{}holder\textgreater{}/root} does not resolve. A
path cannot cross that gap, so a mount does. The holder, the one
namespace that has the mount, clones it detached and movable with
\texttt{open\_tree}, and \texttt{kenneld} relays the descriptor to the
kennel's init over a datagram the init binds at a fixed point in its
view. The init moves the clone into place at \texttt{/dev/binderfs-mesh}
and adds the device to the workload's filesystem ruleset before it forks
the workload, so the broker and facade that use the bus open it by an
ordinary path and know nothing of how it arrived. The clone shares the
original's superblock, so a transaction on it crosses the user-namespace
boundary intact, the property the whole arrangement rests on and the one
most worth having measured rather than reasoned.

\hypertarget{the-brokered-shapes}{%
\section{14.7 The brokered shapes}\label{the-brokered-shapes}}

On the mesh bus the two declarable shapes deliver as their grammar
promised. A \texttt{binder-connector} consume returns the provider's
node handle, which the consumer reaches by name with whatever rides it
opaque to the daemon. A \texttt{dbus-name} consume mints no socket: on
the per-kennel bus the daemon answers the consumer's locator with the
mesh device's path, and the consumer's existing D-Bus facade opens the
mesh bus and reaches its broker there, so a brokered name widens an
allow-set the facade already had rather than opening a new route.

The broker that name reaches is itself a kennel. A standing service
kennel runs \texttt{dbus-broker}, holding one control node on the D-Bus
mesh bus that only \texttt{kenneld} addresses, and the single verb on it
is \texttt{ACCEPT\_SESSION}: the daemon tells the broker to admit one
session under a named filter, the broker mints a fresh per-session node,
stores the filter against that node's cookie, and returns the node.
\texttt{kenneld} forwards it to the consumer, which then speaks to the
broker directly on its own session node, send and receive and close,
with the daemon out of the byte path as everywhere else. When the
consumer drops the session node its release tells the broker to reclaim
the session. The daemon relays no frame; it authorises a session and
steps aside.

What the broker enforces against is the consumer's identity, and the
daemon resolves that without holding a table. On each connect it reads
the kernel-attested pid of the caller, resolves that pid's cgroup to the
kennel context the cgroup names, and looks up that context's one settled
D-Bus filter. The cgroup is the restart-invariant name of the kennel, so
a facade that restarts under a fresh pid resolves to the same identity,
and the only standing state is the policy held against the context and
dropped when the context is released. There is no session cookie to
forge and no per-connection identity to leave stale.

The mesh bus is where the cross-kennel relay residual concentrates, and
delivering the shapes does not dissolve it. A legitimately granted edge
is a channel like any other, so a provider a consumer is entitled to
reach can answer across that reach, and a compromised provider can use
the edge it was granted (T5.2). The grant is the consumer's to make and
the daemon brokers rather than relays, which holds the residual to edges
an operator signed, but a signed edge is real authority once given. The
declaration surface, the catalogue, the matching, the supervision, and
the enablement treat all three shapes alike, so the mesh shapes are the
af-unix shape's siblings and not a second grammar
(\texttt{make-invalid-unrepresentable}); what differs is only the
transport the daemon brokers, never the way a consumer asks.
